\section{Introduction} \label{sec:introduction}
Linux and its command line remain important in computing education and professional practice.
Practical command-line skills continue to be relevant for topics such as system administration, software development, and computer security.
Prior research has shown that command-line interactions support hands-on learning and that terminal skills can be taught through interactive assignments~\cite{svabensky_toolset_2021,bailey_uassign_2019}.
At the University of Tartu, Linux command-line skills play an important role in courses such as ``Operating Systems'' and ``Computer Security''.
Consequently, a web-based Linux learning environment offers educational value by supporting hands-on command-line practice.

Students at the University of Tartu have previously explored this idea in two earlier theses.
Halapuu~\cite{halapuu_2022} developed the initial proof-of-concept web-based Linux command-line learning environment, and Eistre~\cite{eistre_2024} extended it into a cloud-hosted learning platform designed for continuous use.
These earlier iterations demonstrated the feasibility and educational value of the approach.
However, they also showed issues related to maintainability, deployment reproducibility, portability across environments, and the cost efficiency of the deployment model.
In addition, parts of the technical stack and operational model had become outdated, limiting the platform's suitability for long-term use.

The aim of this thesis is to modernize the existing cloud-hosted learning environment while preserving its educational purpose and core workflow.
To achieve this, the thesis introduces clearer architectural modularization, a provider-agnostic provisioning layer for user-specific Linux environments, dedicated scheduled jobs for periodic maintenance, a modernized authentication layer with support for single sign-on, and improved multilingual content management.
The thesis also adopts a reproducible deployment approach using Helm\footnote{\url{https://helm.sh/}} and Terraform\footnote{\url{https://developer.hashicorp.com/terraform}}.
The revised platform is then deployed and evaluated on UT HPC\footnote{\url{https://hpc.ut.ee/}}, Microsoft Azure\footnote{\url{https://azure.microsoft.com/}}, and Amazon Web Services\footnote{\url{https://aws.amazon.com/}}.
The thesis therefore presents a redesigned implementation of the learning environment together with a comparative deployment evaluation across these three target environments.

The thesis is structured as follows.
Chapter~\ref{sec:background-and-previous-system} presents the educational and technical background, describes the two earlier iterations of the learning environment, and identifies the main limitations that motivated the thesis work.
Chapter~\ref{sec:revised-architecture-and-implementation} describes the revised architecture, the main implementation changes, and the development workflow followed during implementation.
Chapter~\ref{sec:deployment-environments-and-infrastructure-setup} explains how the redesigned system was deployed on UT HPC, Microsoft Azure, and Amazon Web Services.
Chapter~\ref{sec:evaluation-and-discussion} evaluates the deployed learning environment during the observed course period and discusses the resulting operational and cost trade-offs.
Chapter~\ref{sec:conclusion} presents the conclusion.
Appendix~\ref{app:deployment-and-maintenance-instructions} provides brief deployment and maintenance instructions.
The remaining appendices provide the supporting evidence used in the cost comparison, including Azure and AWS billing screenshots and UT HPC monitoring material.

For formatting and language refinement, OpenAI's GPT-5.4~\cite{openai_gpt_5_4_2026} was used to identify grammatical issues, point out formatting problems, and suggest clearer phrasing, as well as improvements to the structure and flow of existing text.
AI-based coding assistants were also used during the implementation work described in Chapter~\ref{sec:revised-architecture-and-implementation}, where they supported refactoring planning, code generation, testing, and debugging.
All generated suggestions and code were reviewed, adapted, tested, and integrated by the author.
